
\section{可积性}\label{sec:E-3}

\BookRecordIndex{T0041}Newlander--Nirenberg 定理断言，\BookRecordIndex{T0042}Nijenhuis 张量为零的\BookRecordIndex{T0088}殆复结构都是可积的。在二维，\BookRecordIndex{T0042}Nijenhuis 张量总为零，所以每个\BookRecordIndex{T0088}殆复结构都可积。这等价于说黎曼曲面上的每个度量都共形平坦。下面对 $C^{1+\alpha}$ 类的\BookRecordIndex{T0088}殆复结构和度量证明这些经典结果。

记 $\mathcal J\subset\mathbb R^{2\times2}$ 为满足 $j^2=-\mathbf1$、且对每个非零 $v\in\mathbb R^2$，$\{v,jv\}$ 是 $\mathbb R^2$ 的正定向基的矩阵 $j$ 的集合。以 $i\in\mathcal J$ 记 $\mathbb R^2\cong\mathbb C$ 上的标准复结构。圆盘 $\mathbb D\subset\mathbb C$ 上的坐标写作 $s+it$。

\begin{theorem}[Newlander--Nirenberg]\label{thm:E-3-1}
设 $0<\alpha<1$，$j:\mathbb D\to\mathcal J$ 是 $C^{1+\alpha}$ 类\BookRecordIndex{T0088}殆复结构。则存在局部 $C^{2+\alpha}$ 微分同胚 $\psi:(\mathbb D,0)\to(\mathbb D,0)$，使 $\psi^*i=j$。若 $j$ 光滑，则 $\psi$ 也光滑。若 $j(0)=i$，还可取 $d\psi(0)=\mathbf1$。
\end{theorem}

方程 $\psi^*i=j$ 等价于其逆映射满足 $(\psi^{-1})^*j=i$。因此我们寻找局部 $j$-全纯曲线 $\psi^{-1}:(\mathbb C,i)\to(\mathbb D,j)$。附录 C 发展了求解\BookRecordIndex{T0209}柯西--黎曼方程的技术；定理~\ref{thm:E-3-1} 的证明将变形 $\psi^*i=j$，以便使用这些结果。在给出细节之前，先把定理改述成 E.4 节所需的形式。

$\mathbb D$ 上的黎曼度量可写作
\[
  g=E\,ds^2+2F\,ds\,dt+G\,dt^2,
\]
其中 $E,F,G:\mathbb D\to\mathbb R$ 视情形为光滑或 $C^{1+\alpha}$ 函数，并满足 $E>0$、$G>0$、$EG-F^2>0$。度量 $g$ 给出体积形式 $\mathrm{dvol}=\mathrm{dvol}_g\in\Omega^2(\mathbb D)$ 和\BookRecordIndex{T0088}殆复结构 $j=j_g:\mathbb D\to\mathcal J$：
\[
  \mathrm{dvol}_g:=\sqrt{EG-F^2}\,ds\wedge dt,
  \qquad
  j_g:=\frac1{\sqrt{EG-F^2}}
  \begin{pmatrix}-F&-G\\ E&F\end{pmatrix}.
\]
反过来，由均给出 $\mathbb D$ 标准定向的体积形式与\BookRecordIndex{T0088}殆复结构，可按通常公式 $g=\mathrm{dvol}(\cdot,j\cdot)$ 恢复度量。注意 $j=i$ 当且仅当 $E=G=: \lambda$ 且 $F=0$，即 $g=\lambda(ds^2+dt^2)$ 共形平坦。因此定理~\ref{thm:E-3-1} 可改述如下。

\begin{corollary}[\BookRecordIndex{T0098}等温坐标]\label{cor:E-3-2}
设 $0<\alpha<1$，$g=E\,ds^2+2F\,ds\,dt+G\,dt^2$ 是 $\mathbb D$ 上的 $C^{1+\alpha}$ 类度量。则存在保定向的局部 $C^{2+\alpha}$ 微分同胚 $\psi:(\mathbb D,0)\to(\mathbb D,0)$，使
\[
  \psi^*(ds^2+dt^2)=\lambda(E\,ds^2+2F\,ds\,dt+G\,dt^2),
  \qquad \lambda:=\frac{\det(d\psi)}{\sqrt{EG-F^2}},
\]
或等价地，
\[
  |\partial_s\psi|^2=\lambda E,\qquad
  \langle\partial_s\psi,\partial_t\psi\rangle=\lambda F,\qquad
  |\partial_t\psi|^2=\lambda G.
\]
若 $g$ 光滑，则 $\psi$ 也光滑。若 $E(0)=G(0)$ 且 $F(0)=0$，还可取 $d\psi(0)=\mathbf1$。
\end{corollary}

以 $\Delta_g$ 记度量 $g$ 的 Laplace--Beltrami 算子：
\begin{align*}
\Delta_g:={}&-d^*d\\
={}&\frac{G\partial_s^2-2F\partial_s\partial_t+E\partial_t^2}{EG-F^2}\\
&+\frac1{\sqrt{EG-F^2}}
 \left(\partial_s\frac{G}{\sqrt{EG-F^2}}-\partial_t\frac{F}{\sqrt{EG-F^2}}\right)\partial_s\\
&+\frac1{\sqrt{EG-F^2}}
 \left(\partial_t\frac{E}{\sqrt{EG-F^2}}-\partial_s\frac{F}{\sqrt{EG-F^2}}\right)\partial_t.
\end{align*}
若 $g=ds^2+dt^2$ 为标准度量，则 $\Delta_g=\Delta=\partial_s^2+\partial_t^2$ 为通常的\BookRecordIndex{T0220}拉普拉斯算子；若 $\psi:\mathbb D\to\mathbb D$ 是微分同胚，则
\begin{align*}
\Delta_{\psi^*(ds^2+dt^2)}f
={}&\bigl(\Delta(f\circ\psi^{-1})\bigr)\circ\psi\\
={}&\frac{|\partial_t\psi|^2\partial_s^2f
 -2\langle\partial_s\psi,\partial_t\psi\rangle\partial_s\partial_tf
 +|\partial_s\psi|^2\partial_t^2f}{\det(d\psi)^2}\\
&+\frac1{\det(d\psi)}\left(
 \partial_s\frac{|\partial_t\psi|^2}{\det(d\psi)}
 -\partial_t\frac{\langle\partial_s\psi,\partial_t\psi\rangle}{\det(d\psi)}\right)\partial_sf\\
&+\frac1{\det(d\psi)}\left(
 \partial_t\frac{|\partial_s\psi|^2}{\det(d\psi)}
 -\partial_s\frac{\langle\partial_s\psi,\partial_t\psi\rangle}{\det(d\psi)}\right)\partial_tf.
\end{align*}
这一算子的二阶部分记作
\begin{equation}\label{eq:E-3-1}
  (\psi^*\Delta):=
  \frac{|\partial_t\psi|^2\partial_s^2
  -2\langle\partial_s\psi,\partial_t\psi\rangle\partial_s\partial_t
  +|\partial_s\psi|^2\partial_t^2}{\det(d\psi)^2}.
\end{equation}
据此，推论~\ref{cor:E-3-2} 可再改述如下。

\begin{corollary}\label{cor:E-3-3}
设 $0<\alpha<1$，$a,b,c:\mathbb D\to\mathbb R$ 是满足 $a>0$、$c>0$、$ac-b^2>0$ 的 $C^{1+\alpha}$ 函数。则存在局部 $C^{2+\alpha}$ 微分同胚 $\psi:(\mathbb D,0)\to(\mathbb D,0)$，使
\[
  \psi^*\Delta=\lambda(a\partial_s^2+2b\partial_s\partial_t+c\partial_t^2),
  \qquad \lambda:=\frac1{\det(d\psi)\sqrt{ac-b^2}}.
\]
若 $a,b,c$ 光滑，则 $\psi$ 也光滑。若 $a(0)=c(0)$ 且 $b(0)=0$，还可取 $d\psi(0)=\mathbf1$。
\end{corollary}

\begin{proof}[定理~\ref{thm:E-3-1} 的证明]
全文对任意 $C^1$ 映射 $f:\mathbb D\to\mathbb C$ 记
\begin{equation}\label{eq:E-3-2}
  \partial f:=\frac12(\partial_sf-i\partial_tf),\qquad
  \bar\partial f:=\frac12(\partial_sf+i\partial_tf).
\end{equation}
证明分为五步。先将问题化为求解 Cauchy--Beltrami 方程，再求解该方程。

\emph{第一步。}设 $j:\mathbb D\to\mathcal J$ 是形如
\[
  j=\begin{pmatrix}b&-a\\c&-b\end{pmatrix},
  \qquad a>0,\quad c>0,\quad ac-b^2=1
\]
的 $C^1$ 类\BookRecordIndex{T0088}殆复结构，定义 $\mu:\mathbb D\to\mathbb C$ 为
\[
  \mu:=\frac{c-a-2ib}{a+c+2},
\]
并设 $\psi:\mathbb D\to\mathbb C$ 为 $C^2$ 函数。下列两条等价：
\begin{enumerate}
\item[(i)] $\psi$ 在零点某邻域上的限制是到零点开邻域的微分同胚，且 $\psi^*i=j$、$\psi(0)=0$、$\partial_s\psi(0)=1$；
\item[(ii)] $\psi$ 在零点某邻域满足 Cauchy--Beltrami 方程
\begin{equation}\label{eq:E-3-3}
  \bar\partial\psi=\mu\partial\psi,\qquad
  \psi(0)=0,\qquad \partial_s\psi(0)=1.
\end{equation}
\end{enumerate}
注意
\[
  d\psi\circ j=(b\partial_s\psi+c\partial_t\psi)\,ds
  -(a\partial_s\psi+b\partial_t\psi)\,dt.
\]
所以 $i\,d\psi=d\psi\circ j$ 当且仅当
\[
  (b-i)\partial_s\psi+c\partial_t\psi=0,
  \qquad a\partial_s\psi+(b+i)\partial_t\psi=0.
\]
因 $ac-b^2=1$，这两式等价；它们又与 $\bar\partial\psi=\mu\partial\psi$ 等价，因为
\[
  \bar\partial\psi-\mu\partial\psi
  =\frac{a+1+ib}{a+c+2}
   \left(\partial_s\psi+\frac{b+i}{a}\partial_t\psi\right).
\]
若 $\bar\partial\psi=\mu\partial\psi$ 且 $\partial_s\psi(0)=1$，则 $\partial_s\psi(0)$ 与 $\partial_t\psi(0)$ 在实数域上线性无关。由逆函数定理，$\psi$ 在零点附近是局部微分同胚。第一步得证。

现在设 $j:\mathbb D\to\mathcal J$ 为 $C^{1+\alpha}$ 类\BookRecordIndex{T0088}殆复结构，$\mu:\mathbb D\to\mathbb C$ 为第一步定义的 $C^{1+\alpha}$ 函数。处处有 $|\mu|<1$。下面四步证明 (E.3.3) 在零点附近有局部 $C^{2+\alpha}$ 解 $\psi:(\mathbb D,0)\to\mathbb C$，且当 $\mu$ 光滑时解也光滑。先处理光滑性。

\emph{第二步。}设 $\psi:\mathbb D\to\mathbb C$ 是式~\eqref{eq:E-3-3} 的一个 $C^{1+\alpha}$ 解，且 $W\subset\mathbb D$ 是零点的开邻域，$\psi|_W$ 是到开集 $U\subset\mathbb D$ 的微分同胚。若 $\mu$ 属于 $C^{k+\alpha}$（整数 $k>0$），则 $\psi$ 在 $W$ 上属于 $C^{k+1+\alpha}$。

令 $j:\mathbb D\to\mathcal J$ 为 $\mu$ 对应的\BookRecordIndex{T0088}殆复结构，$u:U\to W$ 为 $\psi$ 的逆映射。于是 $j\in C^{k+\alpha}$，由第一步有 $u^*j=i$，等价地，$\bar\partial_j(u)=0$。特别地，对每个 $p>2$ 有 $j\in W^{k,p}$；依定理 B.4.1 与注 B.4.3，对每个 $p>2$ 有 $u\in W^{k+1,p}$。由定理 B.1.11 \BookRecordIndex{T0326}推出 $u\in C^{k+\alpha}$。此外，$u$ 满足
\[
  \Delta u=\bigl(\partial_t(j\circ u)\bigr)\partial_su
           -\bigl(\partial_s(j\circ u)\bigr)\partial_tu.
\]
由于 $u$ 与 $j$ 局部属于 $C^{k+\alpha}$，故 $\Delta u\in C^{k-1+\alpha}$；再由 Schauder \BookRecordIndex{T0408}正则性（见 Gilbarg--Trudinger~\srcite{140}）得 $u\in C^{k+1+\alpha}$。因此 $\psi\in C^{k+1+\alpha}$，如所称。

\emph{第三步。}若 $\mu(0)=0$，则式~\eqref{eq:E-3-3} 在零点附近有局部 $C^{2+\alpha}$ 解 $\psi$。

对 $\beta>0$ 且 $\beta\notin\mathbb Z$，记 $C^\beta(\mathbb D)=C^\beta(\mathbb D,\mathbb C)$ 为复值函数的相应 Hölder 空间。由定理 C.4.1 和 Schauder \BookRecordIndex{T0408}正则性（见文献~\srcite{140}），\BookRecordIndex{T0210}柯西--黎曼算子 $\bar\partial:C^{1+\alpha}(\mathbb D)\to C^\alpha(\mathbb D)$ 限制在 Banach 空间
\[
  C_0:=\{\xi\in C^{1+\alpha}(\mathbb D)\mid
  \xi(e^{i\theta})\in\mathbb R e^{3i\theta/2}\}
\]
上，给出到 $C^\alpha(\mathbb D)$ 的满射 Fredholm 算子，其实指标为四。核由所有形如
\[
  \xi(z)=a_0+a_1z+\bar a_1z^2+\bar a_0z^3
\]
的多项式构成。因此线性映射
\[
  C_0\longrightarrow C^\alpha(\mathbb D)\times\mathbb C\times\mathbb C,
  \qquad \xi\longmapsto(\bar\partial\xi,\xi(0),\partial_s\xi(0))
\]
是 Banach 空间同构。以 $T:C^\alpha(\mathbb D)\times\mathbb C\times\mathbb C\to C^{1+\alpha}(\mathbb D)$ 记其逆。于是 $T(0,0,1)(z)=z+z^2$；若 $f\in C^\alpha(\mathbb D)$，则 $\psi:=T(f,0,1)$ 满足
\[
  \bar\partial\psi=f,\qquad \psi(0)=0,\qquad \partial_s\psi(0)=1.
\]
特别地，$\psi$ 满足式~\eqref{eq:E-3-3} 的归一化条件。（我们将上述算子的指标设为四，正是为了容纳这些条件。）现定义光滑非线性映射
\[
  \Phi:C^\alpha(\mathbb D)\times C^\alpha(\mathbb D)
  \longrightarrow C^\alpha(\mathbb D)\times C^\alpha(\mathbb D),
  \quad \Phi(f,\lambda):=(f-\lambda\partial(T(f,0,1)),\lambda).
\]
若 $\Phi(f,\lambda)=(0,\lambda)$，则 $\psi:=T(f,0,1)$ 满足 $\bar\partial\psi=f=\lambda\partial\psi$ 及 $\psi(0)=0$、$\partial_s\psi(0)=1$，即以 $\mu=\lambda$ 代入式~\eqref{eq:E-3-3} 所得的三个条件。因此只需证明当 $\lambda(0)=0$ 时，方程 $\Phi(f,\lambda)=(0,\lambda)$ 有局部解。

首先注意
\[
  (d\Phi(0,0)(\hat f,\hat\lambda))(z)
  =(\hat f-\hat\lambda(1+2z),\hat\lambda).
\]
此算子可逆，故由逆函数定理，$\Phi$ 在零点的某个开邻域 $U^\alpha\subset C^\alpha(\mathbb D)\times C^\alpha(\mathbb D)$ 上有光滑逆映射，并将零点映到零点。取光滑\BookRecordIndex{T0188}截断函数 $\chi:\mathbb D\to[0,1]$，在零点附近等于一，在 $|z|\ge1/2$ 时等于零，并在圆盘外延拓为零。\sourcefootnote{原书实体页 682 只将 $\chi$ 定义在 $\mathbb D$ 上，随即对全部 $z\in\mathbb D$ 使用 $\chi(z/\varepsilon)$；当 $|z|>\varepsilon$ 时后者的自变量不一定在 $\mathbb D$。因 $\chi$ 在 $|z|\ge1/2$ 为零，按零延拓即可使该式有定义。}对 $0<\varepsilon\le1$，定义 $\mu_\varepsilon:\mathbb D\to\mathbb C$ 为
\[
  \mu_\varepsilon(z):=\mu(z)\chi(z/\varepsilon).
\]
由 $\mu(0)=0$ 且 $\mu$ 连续可微，当 $\varepsilon\to0$ 时，$\mu_\varepsilon$ 在 $C^\alpha$ 范数下趋于零。因此充分小的 $\varepsilon$ 满足 $(0,\mu_\varepsilon)\in U^\alpha$。于是
\[
  \psi:=T(f,0,1),\qquad (f,\mu_\varepsilon):=\Phi^{-1}(0,\mu_\varepsilon)
\]
是式~\eqref{eq:E-3-3} 的局部 $C^{1+\alpha}$ 解。由第二步，$\psi\in C^{2+\alpha}$。

\emph{第四步。}若 $\mu$ 为常数，则式~\eqref{eq:E-3-3} 在 $\mathbb D$ 上有光滑解。显式公式为
\[
  \psi(s+it):=s+i\frac{1-\mu}{1+\mu}\,t.
\]

\emph{第五步。}证明定理。设 $\mu:\mathbb D\to\mathbb C$ 为处处满足 $|\mu|<1$ 的 $C^{1+\alpha}$ 函数。令 $\varphi:\mathbb D\to\mathbb C$ 是第四步构造的方程
\[
  \bar\partial\varphi=\mu(0)\partial\varphi,\qquad
  \varphi(0)=0,\qquad\partial_s\varphi(0)=1
\]
的解，置
\[
  \widetilde\mu:=
  \left(\frac{\partial\varphi}{\bar\partial\bar\varphi}
  \frac{\mu-\mu(0)}{1-\overline{\mu(0)}\mu}\right)\circ\varphi^{-1}.
\]
因 $\widetilde\mu(0)=0$，第三步保证以 $\widetilde\mu$ 代 $\mu$ 的式~\eqref{eq:E-3-3} 有局部 $C^{2+\alpha}$ 解 $\widetilde\psi:(\mathbb D,0)\to(\mathbb D,0)$。定义 $\psi:=\widetilde\psi\circ\varphi$。由恒等式
\[
  \mu_{\psi\circ\varphi^{-1}}
  =\left(\frac{\partial\varphi}{\bar\partial\bar\varphi}
  \frac{\mu_\psi-\mu_\varphi}
  {1-\bar\mu_\varphi\mu_\psi}\right)\circ\varphi^{-1},
  \qquad \mu_f:=\frac{\bar\partial f}{\partial f},
\]
可知 $\psi$ 是式~\eqref{eq:E-3-3} 的 $C^{2+\alpha}$ 解。若 $\mu$ 光滑，则 $\widetilde\mu$ 光滑，依第二步 $\widetilde\psi$ 光滑，从而 $\psi$ 光滑。第一步表明 $\psi$ 满足定理要求。
\end{proof}
